\documentclass[10pt,a4paper]{amsart}
\usepackage[margin=19mm]{geometry}
\usepackage{amsmath,amssymb,mathtools}
\usepackage[scr=rsfs,cal=euler]{mathalfa}
\usepackage{xfrac}
\usepackage[unicode,hidelinks,bookmarks=false]{hyperref}
\newcommand{\ie}{i.e.\ }
\newcommand{\cf}{cf.\ }
\newcommand{\resp}{respectively\ }
\newcommand{\Subsection}{\subsection}
\providecommand{\nobreakdash}{\nobreak\hbox{-}}
\setlength{\emergencystretch}{2em}
\multlinegap0pt
\usepackage{bm}
\usepackage{bbm}
\usepackage{dsfont}
\usepackage{xspace}
\usepackage{cancel}
\usepackage{mathtools}
\usepackage{amsthm}
\makeatletter
\@ifundefined{theorem}{\newtheorem{theorem}{Theorem}}{}
\@ifundefined{lemma}{\newtheorem{lemma}[theorem]{Lemma}}{}
\makeatother
\title[Ours go to 211: Euler pseudoprimes to 47 prime bases (from C]{Ours go to 211: Euler pseudoprimes to 47 prime bases (from Carmichael numbers)}
\author{Alejandra Alcantarilla S\'anchez, Jolijn Cottaar, Tanja Lange, Benne de Weger}
\date{Source: arXiv:2602.21840v1 (CC BY 4.0)}
\begin{document}
\maketitle

\noindent\emph{Source and licence.} Ours go to 211: Euler pseudoprimes to 47 prime bases (from Carmichael numbers). Version arXiv:2602.21840v1 (CC BY 4.0), distributed under the Creative Commons Attribution 4.0 International licence, as recorded in the arXiv licence metadata for this exact version and shown on the abstract page https://arxiv.org/abs/2602.21840v1. Full rights notice: licenses/arxiv-2602.21840-CC-BY-4.0.txt.\par\smallskip

\noindent\textit{Editorial scope.} Counted unit: the Lemma whose source label is \texttt{le:classification-products} in the article (source0.tex lines 912--948, source label \texttt{le:classification-products}), delivered together with the article's own complete original proof of it (source0.tex lines 949--1218, one \texttt{proof} environment of 270 lines). No mathematics is written, repaired or re-derived: statement and proof are the article's own text line for line. Required context retained uncounted: none. Every definition, display and label of those lines is retained unchanged. Omitted from this package and not redistributed: 1--911, 1219--4136. Nothing inside a retained excerpt is omitted or altered: every retained body is delivered exactly as the article's lines are written, so no omission receipt of any kind accompanies this package. Citations: the retained excerpts carry no citation command at all, so no bibliography entry of the article is transcribed and no reference is redistributed. Reference adaptation: none was needed and none was made; every reference inside the delivered unit resolves inside it, so no unresolved reference or citation is produced. Printed slips kept literal: no printed word, symbol, hypothesis or number of the counted statement or of its proof was altered. Layout and class: the article's own class is \texttt{amsart} with packages including graphicx, amsmath, amsfonts, xcolor, subcaption, xspace, amssymb, algorithm, algpseudocode, mathtools, array, ragged2e, float, setspace; the wrapper is the lane's pinned \texttt{amsart} editorial wrapper at 10pt on A4, which restates the theorem-like environments the retained text needs and adds the article's own macro definitions that the retained text needs (none), with extra wrapper packages (bm, bbm, dsfont, xspace, cancel, mathtools). The delivered numbering is the unit's own. Assets: no retained span contains a figure environment, a graphics-inclusion command, a PDF-page inclusion, a table or a TikZ picture; no image file is copied into this package, the package declares no asset dependency and no image file is redistributed with it. Licence: version arXiv:2602.21840v1 of this article is distributed under the Creative Commons Attribution 4.0 International licence, as recorded in the arXiv licence metadata for this exact version and shown on the abstract page https://arxiv.org/abs/2602.21840v1. A full rights notice is delivered at licenses/arxiv-2602.21840-CC-BY-4.0.txt. Characters: every retained excerpt is pure ASCII, so the delivered engine, XeLaTeX with its default Latin Modern text font, reports no missing character.
\begin{lemma}\label{le:classification-products}
Let
$n_1$
and
$n_2$
with
$\gcd(n_1,n_2)=1$
be Carmichael numbers.
The following are necessary conditions and impossibility results for $n =
n_1
n_2$
to be a Carmichael number:
\begin{itemize}
\item If $n_1$ and $n_2$ are both in class A and $n$ is a Carmichael number, then
\begin{itemize}
\item $n$ is in class A if $\min\{v_2(n_1 - 1), v_2(n_2-1)\} > \max\{v_2(\lambda(n_1)), v_2(\lambda(n_2))\}$;
\item $n$ is in class B1 if  $ \min\{v_2(n_1 - 1), v_2(n_2-1)\} = \max\{v_2(\lambda(n_1)), v_2(\lambda(n_2))\}$ and $v_2(n_1 - 1) \neq v_2(n_2-1)$.
\end{itemize}
Otherwise $n$ is not a Carmichael number.
\item If $n_1$ is in class A, $n_2$ is in class B, and $n$ is a Carmichael number, then
\begin{itemize}
\item $n$ is in class A if $v_2(n_1 - 1) = v_2(n_2-1)$;
\item $n$ is in class B1 if $v_2(n_1 - 1) > v_2(n_2-1)$ and
$v_2(\lambda(n_1))<
v_2(\lambda(n_2))$;
\item $n$ is in class B1 or B2 if $v_2(n_1 - 1) > v_2(n_2-1)$ and
$v_2(\lambda(n_1))=
v_2(\lambda(n_2))$.
\end{itemize}
Otherwise $n$ is not a Carmichael number.
\item If $n_1$ and $n_2$ are both in class B and $n$ is a Carmichael number, then
\begin{itemize}
\item $n$ is in class A if $v_2(n_1 - 1) = v_2(n_2-1)$.
\end{itemize}
Otherwise $n$ is not a Carmichael number.
\end{itemize}
\end{lemma}

\begin{proof}
The proofs for the different cases use that
$v_2(n
-1) =
v_2(n_1n_2
- 1) =
v_2(n_2(n_1
-1) +
n_2
- 1)
\geq
\min
\{v_2(n_1
- 1),
v_2(n_2-1)\}$.
Note that it is larger than the minimum if
$v_2(n_1
- 1) =
v_2(n_2-1)$.
If
$n=n_1
n_2$
is a Carmichael number then
$v_2(\lambda(n))
=
 \max\{v_2(\lambda(n_1)),
v_2(\lambda(n_2))\}$.

{\bf
Let
$n_1$
and
$n_2$
be in class
A.}
By definition,
their indices are even,
i.e.,
$v_2(n_1
-1)
>
v_2(\lambda(n_1))$
and
$v_2(n_2-1)
>
v_2(\lambda(n_2))$.
Let
$v_2(\lambda(n_1))
\geq
v_2(\lambda(n_2))$.

If
$v_2(n_1
-
1)\le
v_2(n_2-1)$
we have that
$v_2(n-1)\ge
\min
\{v_2(n_1
- 1),
v_2(n_2-1)\}
=
v_2(n_1
- 1)
>
v_2(\lambda(n_1))
=
 \max\{v_2(\lambda(n_1)),
v_2(\lambda(n_2))\}=v_2(\lambda(n))$.
Hence $n$ is in class A.
This case is included in
$\min\{v_2(n_1
- 1),
v_2(n_2-1)\}
>
\max\{v_2(\lambda(n_1)),
v_2(\lambda(n_2))\}$.

If
$v_2(n_1
- 1)
>
v_2(n_2-1)$
we have that
$v_2(n-1)
=
\min
\{v_2(n_1
- 1),
v_2(n_2-1)\}
=
v_2(n_2-1)
$.
This is larger than
$v_2(\lambda(n))$,
i.e.
$n$ is in class A,
iff
$\min\{v_2(n_1
- 1),
v_2(n_2-1)\}
>
\max\{v_2(\lambda(n_1)),
v_2(\lambda(n_2))\}$
and equals
$v_2(\lambda(n))$,
i.e.,
$n$ is in class B,
iff
$\min\{v_2(n_1
- 1),
v_2(n_2-1)\}
=
\max\{v_2(\lambda(n_1)),
v_2(\lambda(n_2))\}$.
In the latter case
$v_2(n_2-1)
>
v_2(\lambda(n_2))$
implies
$v_2(\lambda(n_1))
>
v_2(\lambda(n_2))$
and thus $n$ is in class B1.

Otherwise
$v_2(n-1)
<
v_2(\lambda(n))$
and $n$ is not a Carmichael number.

{\bf
 Let
$n_1$
be in class A and
$n_2$
be in class
B.}
By definition,
$v_2(n_1
-1)
>
v_2(\lambda(n_1))$
and
$v_2(n_2-1)
=
v_2(\lambda(n_2))$.

If
$v_2(n_1
- 1) =
v_2(n_2-1)$
we have that
$v_2(n-1)>v_2(n_2-1)
=v_2(\lambda(n_2))
=
v_2(\lambda(n))$,
because
$v_2(n_2
- 1) =
v_2(n_1-1)
 >
v_2(\lambda(n_1))$
implies that
$v_2(\lambda(n_1))
<
v_2(\lambda(n_2))
 =
v_2(\lambda(n))$,
and $n$ is in class A.

If
$v_2(n_1
- 1)
>
v_2(n_2-1)$
we have that
$v_2(n-1)
=
v_2(n_2-1)
=v_2(\lambda(n_2))\le
\max\{v_2(\lambda(n_1)),
v_2(\lambda(n_2))\}
=
v_2(\lambda(n))$.
If $n$ is a Carmichael number we must have equality,
i.e.,
$n$ is in class B,
which requires
$v_2(\lambda(n_1))
\le
v_2(\lambda(n_2))$.
If
$v_2(\lambda(n_1))
<
v_2(\lambda(n_2))$
there are at least two prime factors
$p_1|n_1$
and
$p_2|n_2$
with
$v_2(p_1-1)
<
v_2(p_2-1)$
and thus $n$ is in class B1.
If
$v_2(\lambda(n_1))
=
v_2(\lambda(n_2))$
the result is in B2 if all prime factors
$p_i$
of
$n_1$
and
$n_2$
have the same
$v_2(p_i
-1)$ and otherwise in B1.

If
$v_2(n_1
- 1)
<
v_2(n_2-1)$
we have that
$v_2(n-1)=v_2(n_1-1)
<
v_2(n_2-1)
=
v_2(\lambda(n_2))
\le \max\{v_2(\lambda(n_1)), v_2(\lambda(n_2))\} = v_2(\lambda(n))$, contradicting that
$n$ is a Carmichael number.

{\bf
 Let
$n_1$
and
$n_2$
be in class
B.}
By definition,
$v_2(n_1-1)
=
v_2(\lambda(n_1))$
and
$v_2(n_2-1)
=
v_2(\lambda(n_2))$.

If
$v_2(n_1
- 1) =
v_2(n_2-1)$,
we have
$v_2(n-1)
>
v_2(n_1-1)
=
v_2(\lambda(n_1))
=
v_2(\lambda(n))$,
so $n$ is in class A.
In all other cases we have
$v_2(n
-1)
<
v_2(\lambda(n))$,
contradicting that $n$ is a Carmichael number.
\end{proof}

\end{document}
